% !TEX root = ../main.tex
%% Devoir surveillé n° 6 — chapitres 17 à 20
%% Résultats vérifiés par verifications/eval_ds.py.

\sujetEval{DS6}{MPSI · Devoir surveillé}{Devoir surveillé \numero6}
  {Durée : \textbf{4 heures}\\Calculatrice interdite\\Chapitres 17 à 20}{4}
  {Exercice 1/L'urne de Pólya/6 pts,Exercice 2/Produit scalaire/5 pts,%
   Exercice 3/Séries numériques/5 pts,Exercice 4/Deux variables/4 pts}
\consigne{Les quatre exercices sont indépendants. On rappelle que
$\sum_{n\geqslant0}\frac{1}{n!}$ converge et a pour somme $\mathrm{e}$.}

\begin{exoEval}{6 points}{L'urne de Pólya}
Une urne contient initialement une boule rouge et une boule blanche. À
chaque étape, on tire une boule au hasard, on note sa couleur, puis on la
remet dans l'urne \emph{avec une boule supplémentaire de la même couleur}.
On note $X_n$ le nombre de boules rouges après $n$ tirages ($X_0=1$) et
$R_k$ l'événement « la $k$-ième boule tirée est rouge ».
\begin{questions}
  \item Déterminer les lois de $X_1$ et de $X_2$.
  \item Montrer par récurrence que $X_n$ suit la loi uniforme sur
    $\ens{1,\dots,n+1}$.
  \item En déduire $E(X_n)$, $V(X_n)$, puis $P(R_{n+1})$.
  \item Calculer $P_{R_2}(R_1)$. Commenter.
\end{questions}
\end{exoEval}

\begin{exoEval}{5 points}{Produit scalaire}
Sur $E=\R_2[X]$, on pose
$\langle P,Q\rangle=P(0)Q(0)+P(1)Q(1)+P(2)Q(2)$.
\begin{questions}
  \item Montrer que $\langle\cdot,\cdot\rangle$ est un produit scalaire
    sur $E$.
  \item Soient $L_0,L_1,L_2$ les polynômes de Lagrange associés à
    $0,1,2$. Montrer que $(L_0,L_1,L_2)$ est une base orthonormée de $E$
    et donner les coordonnées d'un polynôme $P$ dans cette base.
  \item Déterminer le projeté orthogonal de $X^{2}$ sur $F=\R_1[X]$ et la
    distance de $X^{2}$ à $F$.
  \item Interpréter le résultat en termes de droite approchant au mieux
    les points $(0,0)$, $(1,1)$ et $(2,4)$.
\end{questions}
\end{exoEval}

\begin{exoEval}{5 points}{Séries numériques}
\begin{questions}
  \item Déterminer la nature de
    $\sum_{n\geqslant1}\Bigl(\dfrac1n-\ln\Bigl(1+\dfrac1n\Bigr)\Bigr)$.
  \item Trouver $a,b,c\in\R$ tels que $u_n=\dfrac{an^{2}+bn+c}{2^{n}}$
    vérifie $u_n-u_{n+1}=\dfrac{n^{2}}{2^{n}}$ pour tout $n$. En déduire la
    somme $\sum_{n=1}^{+\infty}\frac{n^{2}}{2^{n}}$.
  \item On note $S_N=\sum_{n=2}^{N}\ln\Bigl(1+\dfrac{(-1)^{n}}{n}\Bigr)$.
    Calculer $S_{2N+1}$ et $S_{2N}$. En déduire que la série converge et
    donner sa somme.
  \item Montrer que $\sum_{n\geqslant0}\dfrac{n+1}{n!}$ converge et calculer
    sa somme.
\end{questions}
\end{exoEval}

\begin{exoEval}{4 points}{Deux variables}
Soit $f:(x,y)\mapsto x^{2}+xy+y^{2}-3x-3y$, définie sur $\R^{2}$.
\begin{questions}
  \item Justifier que $f$ est de classe $\mathcal{C}^{1}$, calculer ses
    dérivées partielles et déterminer ses points critiques.
  \item On pose $u=x-1$ et $v=y-1$. Montrer que
    \[f(x,y)=u^{2}+uv+v^{2}-3\qquad\text{et}\qquad
      u^{2}+uv+v^{2}=\Bigl(u+\frac v2\Bigr)^{2}+\frac34\,v^{2}.\]
    Conclure sur les extremums de $f$.
  \item Donner l'équation du plan tangent à la surface $z=f(x,y)$ en
    l'origine, et la direction de plus forte décroissance de $f$ en
    $(0,0)$.
\end{questions}
\end{exoEval}
